\section{Metodología}
% ══════════════════════════════════════════════════════════════════════════════

\subsection{Metodología de desarrollo}

El proyecto adoptará \textbf{Scrum} como marco de trabajo ágil, complementado con
prácticas de \textbf{Ingeniería de Requerimientos} basadas en el estándar IEEE~830
y principios de diseño centrado en el usuario (UCD). Esta combinación permite
gestionar incertidumbre, incorporar retroalimentación continua y entregar valor
incremental a lo largo del ciclo de vida del desarrollo.

La elección de Scrum se fundamenta en el análisis de plataformas de referencia:

\begin{itemize}
  \item \textbf{PLADDES -- UNSAAC} (\texttt{tramite.unsaac.edu.pe}): sistema actual
    de trámite documentario de la UNSAAC que gestiona solicitudes virtuales
    clasificadas por tipo de usuario (estudiante, docente, administrativo, público
    general), seguimiento de expedientes y adjunto de vouchers de pago.
    Constituye el mayor referente directo del proyecto; sus limitaciones de
    usabilidad y ausencia de notificaciones proactivas son el punto de partida
    para las mejoras propuestas.

  \item \textbf{Portal Administrativo -- PUCP}: plataforma que centraliza
    información de trámites paso a paso para toda la comunidad universitaria,
    integra pagos académicos virtuales y ofrece paneles diferenciados por
    dirección. Destaca por su arquitectura orientada a servicios y la
    digitalización progresiva de ventanillas presenciales.

  \item \textbf{SISADES -- Universidad Nacional de Huancavelica}: sistema
    integral que centraliza documentos institucionales, biblioteca virtual y
    seguimiento de egresados. Su mesa de partes virtual y el módulo de
    seguimiento de trámites publicado en la Plataforma del Estado Peruano
    (gob.pe) evidencian una integración con estándares de interoperabilidad
    pública que el presente proyecto también considerará.
\end{itemize}

\subsubsection{Roles Scrum}

\begin{center}
\begin{tabular}{p{3.5cm} p{9.5cm}}
  \toprule
  \rowcolor{unsaacRed}
  \color{white}\textbf{Rol} & \color{white}\textbf{Responsabilidad} \\
  \midrule
  Product Owner   & Define y prioriza el \textit{Product Backlog} en función de los
                    objetivos del TUPA y las necesidades de estudiantes y personal
                    administrativo. \\
  Scrum Master    & Facilita ceremonias, elimina impedimentos y asegura que el equipo
                    siga el marco de trabajo. \\
  Dev Team        & Diseña, implementa y prueba los incrementos de producto en cada
                    sprint. \\
  \bottomrule
\end{tabular}
\end{center}

\subsubsection{Ceremonias y artefactos}

\begin{itemize}
  \item \textbf{Sprint Planning}: primera sesión del sprint; se seleccionan ítems
    del backlog y se define el \textit{Sprint Goal}.
  \item \textbf{Daily Stand-up} (10\,min): sincronización diaria de avances,
    impedimentos y planes del día.
  \item \textbf{Sprint Review}: demostración del incremento al \textit{Product
    Owner} y recogida de retroalimentación.
  \item \textbf{Sprint Retrospective}: mejora continua del proceso.
  \item \textbf{Artefactos}: Product Backlog, Sprint Backlog, Definition of Done
    e Incremento potencialmente entregable.
\end{itemize}

\subsubsection{Stack tecnológico propuesto}

\begin{center}
\begin{tabular}{p{3.8cm} p{4.2cm} p{5cm}}
  \toprule
  \rowcolor{unsaacRed}
  \color{white}\textbf{Capa} & \color{white}\textbf{Tecnología} &
  \color{white}\textbf{Justificación} \\
  \midrule
  Frontend        & React.js + Tailwind CSS  & SPA reactiva, ecosistema maduro,
                                               facilita prototipos interactivos. \\
  Backend         & Node.js + Express.js     & API REST ligera, JSON nativo,
                                               alto rendimiento para I/O. \\
  Base de datos   & PostgreSQL               & Soporte robusto de relaciones,
                                               transacciones y consultas complejas. \\
  Autenticación   & JWT + OAuth2             & Seguridad estándar para sesiones
                                               y roles diferenciados. \\
  Notificaciones  & Nodemailer + WebSockets  & Alertas por correo y en tiempo
                                               real sobre estado de trámites. \\
  Control de versiones & Git + GitHub        & Flujo colaborativo con ramas por
                                               feature y pull requests. \\
  Contenerización & Docker / Docker Compose  & Entornos reproducibles para
                                               desarrollo y despliegue. \\
  \bottomrule
\end{tabular}
\end{center}

% ──────────────────────────────────────────────────────────────────────────────
\subsection{Actividades principales}

El proyecto se organiza en \textbf{cuatro sprints} de dos semanas cada uno, más
una fase inicial de planificación (\textit{Sprint 0}), alineados con las dos
entregas formales del semestre.

\subsubsection{Sprint 0 -- Planificación y análisis de requerimientos\\(hasta 28/05/2026 -- Entrega~1)}

\begin{itemize}
  \item Revisión y análisis de la plataforma PLADDES--UNSAAC vigente.
  \item Estudio comparativo de los sistemas de referencia (PUCP, UNH).
  \item Levantamiento de requerimientos funcionales y no funcionales mediante
    entrevistas con usuarios representativos (estudiantes, administrativos).
  \item Elaboración del \textit{Product Backlog} con historias de usuario
    priorizadas (MoSCoW).
  \item Diseño de la arquitectura lógica y física del sistema.
  \item \textbf{Diseño de la base de datos}: modelo entidad-relación y diagrama
    relacional (PostgreSQL).
  \item \textbf{Prototipado dinámico} de las interfaces principales (Figma):
    pantalla de consulta del TUPA, registro de trámite, seguimiento de
    expediente y panel administrativo.
  \item Configuración del repositorio GitHub con estrategia de ramificación
    (\textit{Git Flow}).
  \item Redacción del Plan de Proyecto (presente documento).
\end{itemize}

\subsubsection{Sprint 1 -- Módulo de consulta TUPA e interfaz base
(01/06/2026 -- 12/06/2026)}

\begin{itemize}
  \item Implementación del módulo de \textbf{consulta de procedimientos}: listado,
    búsqueda por categoría, detalle de requisitos y costos asociados del TUPA.
  \item Desarrollo de la \textbf{autenticación y gestión de usuarios}: registro,
    inicio de sesión, roles (estudiante, administrativo, administrador).
  \item Construcción del \textbf{layout base} de la aplicación: barra de
    navegación, footer, diseño responsivo y sistema de estilos.
  \item Configuración del entorno de base de datos y migraciones iniciales.
  \item Pruebas unitarias de los endpoints REST y pruebas de usabilidad del
    módulo de consulta.
\end{itemize}

\subsubsection{Sprint 2 -- Módulo de registro y seguimiento de trámites
(15/06/2026 -- 26/06/2026)}

\begin{itemize}
  \item Implementación del \textbf{formulario de registro de trámite}: datos del
    solicitante, selección de procedimiento TUPA, adjunto de documentos digitales
    y clave de pago.
  \item Desarrollo del \textbf{módulo de seguimiento de expedientes}: estado
    actual, historial de movimientos y observaciones por oficina.
  \item \textbf{Sistema de notificaciones}: alertas por correo electrónico ante
    cambios de estado (recibido, en revisión, observado, resuelto).
  \item Validación digital de documentos adjuntos (tipo, tamaño, formato PDF).
  \item Pruebas de integración del flujo completo de registro y seguimiento.
\end{itemize}

\subsubsection{Sprint 3 -- Panel administrativo y módulo de reportes
(29/06/2026 -- 10/07/2026)}

\begin{itemize}
  \item Desarrollo del \textbf{panel de control administrativo}: bandeja de
    trámites recibidos, cambio de estado, asignación a oficina responsable y
    registro de observaciones.
  \item Implementación del \textbf{módulo de reportes y estadísticas}: trámites
    procesados por período, tiempo de atención promedio, procedimientos más
    solicitados y exportación a PDF/Excel.
  \item Mejoras de \textbf{accesibilidad} (WCAG 2.1 nivel AA): contraste,
    etiquetas ARIA, navegación por teclado.
  \item Pruebas de carga y rendimiento básicas.
\end{itemize}

\subsubsection{Sprint 4 -- Integración, pruebas finales y documentación
(13/07/2026 -- primera semana de agosto 2026 -- Entrega 2)}

\begin{itemize}
  \item Integración de todos los módulos y pruebas de regresión del sistema.
  \item Corrección de defectos identificados en sprints anteriores.
  \item Despliegue de la versión de demostración en entorno de pruebas (Docker).
  \item Elaboración de la \textbf{documentación técnica}: diagramas de
    arquitectura, manual de usuario, diccionario de datos y guía de instalación.
  \item Actualización del repositorio GitHub con commits organizados por módulo.
  \item Preparación de la \textbf{demostración funcional} de la plataforma
    (Entrega Semestral N.º 2, primera semana de agosto de 2026).
\end{itemize}

% ──────────────────────────────────────────────────────────────────────────────
\subsection{Infraestructura y recursos}

\subsubsection{Hardware}

\begin{center}
\begin{tabular}{p{4.5cm} p{4cm} p{4.5cm}}
  \toprule
  \rowcolor{unsaacRed}
  \color{white}\textbf{Recurso} & \color{white}\textbf{Especificación mínima} &
  \color{white}\textbf{Uso} \\
  \midrule
  Laptops de desarrollo   & 8\,GB RAM, CPU i5/Ryzen 5, SSD 256\,GB &
                            Desarrollo y pruebas locales. \\
  Servidor de despliegue  & VPS 2\,vCPU, 4\,GB RAM, 40\,GB SSD     &
                            Entorno de pruebas y demo. \\
  Red universitaria       & Conexión a internet estable              &
                            Acceso a repositorios y herramientas cloud. \\
  \bottomrule
\end{tabular}
\end{center}

\subsubsection{Software y servicios}

\begin{center}
\begin{tabular}{p{4cm} p{3.5cm} p{5.5cm}}
  \toprule
  \rowcolor{unsaacRed}
  \color{white}\textbf{Herramienta} & \color{white}\textbf{Tipo} &
  \color{white}\textbf{Propósito} \\
  \midrule
  Visual Studio Code      & IDE               & Desarrollo del frontend y backend. \\
  Node.js v20 LTS         & Runtime           & Ejecución del servidor Express. \\
  React.js v18            & Framework UI      & Construcción de la SPA. \\
  PostgreSQL 16           & SGBD              & Persistencia de datos. \\
  Docker Desktop          & Contenerización   & Entornos de desarrollo homogéneos. \\
  GitHub                  & Repositorio       & Control de versiones y CI básico. \\
  Figma                   & Diseño UX/UI      & Prototipado dinámico. \\
  Postman                 & Pruebas API       & Validación de endpoints REST. \\
  Trello / Jira (libre)   & Gestión ágil      & Tablero Scrum y seguimiento de tareas. \\
  Nodemailer + SMTP       & Notificaciones    & Envío de alertas por correo. \\
  \bottomrule
\end{tabular}
\end{center}

\subsubsection{Recursos humanos}

El equipo de desarrollo está conformado por \textbf{4 estudiantes} de la
Escuela Profesional de Ingeniería Informática y de Sistemas de la UNSAAC,
quienes asumirán de manera rotativa los roles de \textit{Product Owner},
\textit{Scrum Master} y miembros del \textit{Development Team}.

Como equipo externo de apoyo se contempla la participación de:
\begin{itemize}
  \item \textbf{Docentes asesores} de la asignatura, en calidad de
    \textit{stakeholders} y revisores de entregables.
  \item \textbf{Personal administrativo de la UNSAAC} (muestra representativa)
    para validación de requerimientos y pruebas de usabilidad.
  \item \textbf{Estudiantes voluntarios} para pruebas funcionales y de
    accesibilidad.
\end{itemize}

% ──────────────────────────────────────────────────────────────────────────────
\subsection{Evaluación y monitoreo}

\subsubsection{Mecanismos de control de calidad}

\begin{itemize}
  \item \textbf{Revisión de código} (\textit{Code Review}): toda funcionalidad
    se integra a la rama principal mediante \textit{pull requests} revisados por
    al menos un compañero de equipo antes de aprobar el merge.
  \item \textbf{Definition of Done (DoD)}: un ítem del backlog se considera
    terminado cuando cuenta con código funcional, pruebas unitarias aprobadas,
    documentación básica actualizada y revisión de pares completada.
  \item \textbf{Pruebas unitarias e integración}: uso de \textit{Jest} para el
    backend y \textit{React Testing Library} para el frontend, con cobertura
    mínima del 70\,\% en módulos críticos.
  \item \textbf{Pruebas de usabilidad}: al menos una sesión de prueba con usuarios
    reales por sprint, utilizando la \textit{Escala de Usabilidad del Sistema}
    (SUS) para medir la experiencia de usuario.
\end{itemize}

\subsubsection{Seguimiento del progreso}

\begin{itemize}
  \item \textbf{Tablero Scrum}: visualización del estado de las tareas
    (\textit{To Do} / \textit{In Progress} / \textit{Done}) actualizado
    diariamente.
  \item \textbf{Burndown Chart}: gráfica de trabajo restante vs. tiempo para
    cada sprint, utilizada en la \textit{Sprint Review} para medir la velocidad
    del equipo.
  \item \textbf{Commits en GitHub}: historial versionado que refleja el avance
    diario; se revisará el número de commits y la cobertura de las historias
    de usuario en cada entrega formal.
  \item \textbf{Bitácora de sprints}: documento que registra impedimentos,
    decisiones técnicas y lecciones aprendidas tras cada retrospectiva.
\end{itemize}

\subsubsection{Indicadores de desempeño (KPIs)}

\begin{center}
\begin{tabular}{p{5.5cm} p{3cm} p{4.5cm}}
  \toprule
  \rowcolor{unsaacRed}
  \color{white}\textbf{Indicador} & \color{white}\textbf{Meta} &
  \color{white}\textbf{Herramienta de medición} \\
  \midrule
  Cobertura de pruebas unitarias   & $\geq$ 70\,\%      & Jest / Coverage report \\
  Puntuación SUS (usabilidad)      & $\geq$ 68 puntos   & Cuestionario SUS \\
  Historias completadas por sprint & $\geq$ 80\,\%      & Burndown chart \\
  Tiempo de respuesta de la API    & $\leq$ 500\,ms     & Postman / logs \\
  Defectos críticos en entrega     & 0                  & Registro de bugs (GitHub Issues) \\
  Commits por sprint               & $\geq$ 20          & GitHub insights \\
  \bottomrule
\end{tabular}
\end{center}

\subsubsection{Cronograma de sprints}

\begin{figure}[H]
\centering
\begin{ganttchart}[
    hgrid,
    vgrid,
    x unit=0.46cm,
    y unit title=0.6cm,
    y unit chart=0.55cm,
    title label font=\scriptsize,
    bar label font=\scriptsize,
    bar/.style={fill=unsaacRed!70},
    bar incomplete/.style={fill=unsaacGray!40},
    milestone/.style={fill=unsaacRed, shape=diamond, inner sep=1.5pt},
    today=10,
    today rule/.style={draw=unsaacRed, ultra thick},
    title height=1
  ]{1}{22}
  \gantttitle{Mayo}{2}
  \gantttitle{Junio}{9}
  \gantttitle{Julio}{9}
  \gantttitle{Ago}{2} \\
  \gantttitle{S4}{1}\gantttitle{S5}{1}
  \gantttitle{S1}{1}\gantttitle{S2}{1}\gantttitle{S3}{1}\gantttitle{S4}{1}
  \gantttitle{S1}{1}\gantttitle{S2}{1}\gantttitle{S3}{1}
  \gantttitle{S1}{1}\gantttitle{S2}{1}\gantttitle{S3}{1}\gantttitle{S4}{1}
  \gantttitle{S5}{1}\gantttitle{S1}{1}\gantttitle{S2}{1}\gantttitle{S3}{1}
  \gantttitle{S4}{1}\gantttitle{S5}{1}
  \gantttitle{S1}{1}\gantttitle{S2}{1}\gantttitle{S3}{1} \\
  % Sprint 0
  \ganttbar{Sprint 0: Análisis y planif.}{1}{2} \\
  \ganttmilestone{Entrega 1 (28/05)}{2} \\
  % Sprint 1
  \ganttbar{Sprint 1: Consulta TUPA}{3}{6} \\
  % Sprint 2
  \ganttbar{Sprint 2: Registro/seguim.}{7}{10} \\
  % Sprint 3
  \ganttbar{Sprint 3: Panel admin.}{11}{14} \\
  % Sprint 4
  \ganttbar{Sprint 4: Integración/doc.}{15}{21} \\
  \ganttmilestone{Entrega 2 (1.ª sem. ago.)}{22}
\end{ganttchart}
\caption{Cronograma de sprints del proyecto (semanas S = semana del mes).}
\end{figure}